\documentclass{article}
\usepackage{amsmath,amssymb,amsthm}
\usepackage{algorithm}
\usepackage{algpseudocode}
\usepackage{fullpage}
\newtheorem{theorem}{Theorem}
\newtheorem{lemma}{Lemma}
\newtheorem{assumption}{Assumption}
\newcommand{\E}{\mathbb{E}}
\newcommand{\R}{\mathbb{R}}

\begin{document}

%=====================================================================
\section*{Differentially Private Stochastic Gradient Descent (DP‑SGD) with Gaussian Noise}
%=====================================================================

\section*{Mathematical Formulation}
Let $f:\R^{d}\rightarrow\R$ be the empirical risk
\[
f(w)=\frac{1}{n}\sum_{i=1}^{n}\ell(w;z_i),
\]
where $\ell$ is convex and $L$–smooth in $w$.  
At iteration $t=0,1,\dots,T-1$ we perform

\[
\boxed{
\begin{aligned}
g_t &:= \nabla f(w_t)+\xi_t,\\
w_{t+1}&:= w_t-\eta_t g_t,
\end{aligned}}
\qquad \xi_t\sim\mathcal N(0,\sigma^{2}I_{d}),
\tag{1}
\]

where $\eta_t>0$ is the stepsize and $\sigma^{2}$ is the variance of the added Gaussian noise.  
The privacy accountant (e.g. moments accountant) updates the cumulative privacy loss
$\varepsilon_t$ after each iteration; the algorithm stops when a pre‑specified budget
$(\varepsilon,\delta)$ is exhausted.

\section*{Pseudocode}
\begin{algorithm}[H]
\caption{DP‑SGD with Gaussian Noise}
\begin{algorithmic}[1]
\Require Initial point $w_0\in\R^{d}$, stepsizes $\{\eta_t\}_{t=0}^{T-1}$, noise scale $\sigma$, privacy budget $(\varepsilon,\delta)$.
\State Initialize privacy accountant $\mathcal A\leftarrow (0,0)$.
\For{$t=0$ \textbf{to} $T-1$}
    \State Compute (full‑batch or minibatch) gradient $\nabla f(w_t)$.
    \State Sample $\xi_t\sim\mathcal N(0,\sigma^{2}I_{d})$.
    \State Form noisy gradient $g_t=\nabla f(w_t)+\xi_t$.
    \State Update $w_{t+1}=w_t-\eta_t g_t$.
    \State Update accountant $\mathcal A$ with the Gaussian mechanism of variance $\sigma^{2}$.
    \If{$\mathcal A$ exceeds $(\varepsilon,\delta)$}
        \State \textbf{break}
    \EndIf
\EndFor
\State \Return $w_T$.
\end{algorithmic}
\end{algorithm}

\section*{Dry Run Example}
Consider the one‑dimensional quadratic
\[
f(w)=(w-3)^{2},\qquad \nabla f(w)=2(w-3).
\]
Set $w_0=0$, constant stepsize $\eta=0.1$, and Gaussian noise variance $\sigma^{2}=0.01$.
For illustration we fix the random draws:
\[
\xi_0=+0.05,\;\xi_1=-0.08,\;\xi_2=+0.02 .
\]

\begin{center}
\begin{tabular}{c|c|c|c|c}
$t$ & $w_t$ & $\nabla f(w_t)$ & $g_t=\nabla f(w_t)+\xi_t$ & $w_{t+1}=w_t-\eta g_t$ \\ \hline
0 & $0$ & $2(0-3)=-6$ & $-6+0.05=-5.95$ & $0-0.1(-5.95)=0.595$ \\[2mm]
1 & $0.595$ & $2(0.595-3)=-4.81$ & $-4.81-0.08=-4.89$ & $0.595-0.1(-4.89)=1.084$ \\[2mm]
2 & $1.084$ & $2(1.084-3)=-3.832$ & $-3.832+0.02=-3.812$ & $1.084-0.1(-3.812)=1.465$ 
\end{tabular}
\end{center}

The objective values after each iteration are
\[
f(w_1)= (0.595-3)^2\approx5.77,\;
f(w_2)= (1.084-3)^2\approx3.68,\;
f(w_3)= (1.465-3)^2\approx2.35 .
\]
The noisy updates move $w_t$ toward the optimum $w^{*}=3$ while respecting a (hypothetical) privacy budget.

%=====================================================================
\section*{Assumptions}
%=====================================================================
\begin{assumption}[Convexity]\label{ass:convex}
$f$ is convex: $\forall w,w'\in\R^{d},\; f(w')\ge f(w)+\langle\nabla f(w),w'-w\rangle$.
\end{assumption}

\begin{assumption}[L‑smoothness]\label{ass:smooth}
There exists $L>0$ such that $\forall w,w'\in\R^{d}$
\[
f(w')\le f(w)+\langle\nabla f(w),w'-w\rangle+\frac{L}{2}\|w'-w\|^{2}.
\]
\end{assumption}

\begin{assumption}[Bounded Gradient]\label{ass:bound}
There exists $G>0$ with $\|\nabla f(w)\|\le G$ for all $w$.
\end{assumption}

\begin{assumption}[Noise Distribution]\label{ass:noise}
The added noise satisfies $\xi_t\sim\mathcal N(0,\sigma^{2}I_d)$ and is independent of $w_t$.
Hence $\E[\xi_t]=0$ and $\E[\|\xi_t\|^{2}]=\sigma^{2}d$.
\end{assumption}

All constants $L,G,\sigma$ are known to the analyst and the stepsizes $\{\eta_t\}$ are deterministic.

%=====================================================================
\section*{Main Convergence Theorem}
%=====================================================================
\begin{theorem}[Expected Suboptimality of DP‑SGD]\label{thm:conv}
Assume \ref{ass:convex}–\ref{ass:noise} and let $\eta_t\equiv\eta\le \frac{1}{L}$.  
Define $w^{*}\in\arg\min_{w}f(w)$. Then for any $T\ge1$,
\[
\E\!\left[f(w_T)-f(w^{*})\right]\;\le\;
\frac{\|w_0-w^{*}\|^{2}}{2\eta T}
\;+\;\frac{\eta}{2}\bigl(G^{2}+\sigma^{2}d\bigr).
\tag{2}
\]
In particular, choosing $\eta=\frac{\|w_0-w^{*}\|}{\sqrt{(G^{2}+\sigma^{2}d)T}}$ yields
\[
\E\!\left[f(w_T)-f(w^{*})\right]=
O\!\left(\frac{1}{\sqrt{T}}\right).
\]
\end{theorem}

%=====================================================================
\section*{Theoretical Properties}
%=====================================================================
\begin{itemize}
\item \textbf{Stability.} The recursion (1) is a contraction in expectation when $\eta\le 1/L$; the added Gaussian term does not destroy stability because its mean is zero and its variance appears only as an additive constant in (2).
\item \textbf{Hyperparameter Sensitivity.}
    \begin{enumerate}
        \item Larger $\eta$ accelerates the $1/(\eta T)$ term but inflates the variance term $\eta(G^{2}+\sigma^{2}d)/2$.
        \item The optimal $\eta$ balances the two terms, giving the $O(1/\sqrt{T})$ rate above.
    \end{enumerate}
\item \textbf{Comparison to Non‑Private SGD.}  Setting $\sigma=0$ recovers the classical bound for (deterministic) SGD on convex $L$‑smooth functions.
\item \textbf{Privacy‑Utility Trade‑off.}  The variance term $\sigma^{2}d$ grows with the privacy noise scale; tighter privacy (smaller $\varepsilon$) forces larger $\sigma$, thus degrading the bound linearly in $\sigma$.
\end{itemize}

\section*{Discussion}
The theorem shows that DP‑SGD retains the same $O(1/\sqrt{T})$ convergence order as standard SGD, with an additional additive penalty proportional to the noise variance required for differential privacy. By carefully calibrating $\sigma$ (via the moments accountant) and the stepsize, one can achieve a desirable balance between privacy guarantees and optimization accuracy.

%=====================================================================
\section*{Preliminaries (Definitions and Lemmas)}
%=====================================================================
\begin{lemma}[Smoothness Inequality]\label{lem:smooth}
If $f$ is $L$‑smooth then for any $w,u\in\R^{d}$
\[
f(u)\le f(w)+\langle\nabla f(w),u-w\rangle+\frac{L}{2}\|u-w\|^{2}.
\]
\end{lemma}
\begin{proof}
Standard from the definition of $L$‑smoothness. \qedhere
\end{proof}

\begin{lemma}[Three‑Point Identity]\label{lem:three}
For any vectors $a,b,c\in\R^{d}$ and scalar $\eta>0$,
\[
\|a-c\|^{2}= \|a-b\|^{2}+2\langle a-b,c-b\rangle+\|c-b\|^{2}.
\]
\end{lemma}
\begin{proof}
Expand $\|a-c\|^{2}=\langle a-c,a-c\rangle$ and collect terms. \qedhere
\end{proof}

%=====================================================================
\section*{Main Proof (Step‑by‑Step Derivation)}
%=====================================================================
We prove Theorem~\ref{thm:conv}.  Throughout the proof expectations are taken w.r.t. the noise $\xi_0,\dots,\xi_{t-1}$ conditioned on the past iterates.

\bigskip
\noindent\textbf{[Step 1]}  Start from the squared distance to the optimum and apply Lemma~\ref{lem:three} with $a=w_{t+1}$, $b=w_t$, $c=w^{*}$:
\[
\|w_{t+1}-w^{*}\|^{2}
= \|w_t-w^{*}\|^{2}
-2\eta\langle g_t,\,w_t-w^{*}\rangle
+\eta^{2}\|g_t\|^{2}.
\tag{3}
\]

\noindent\textbf{[Step 2]}  Decompose $g_t=\nabla f(w_t)+\xi_t$ and rewrite the inner product:
\[
\langle g_t,\,w_t-w^{*}\rangle
= \langle\nabla f(w_t),w_t-w^{*}\rangle
+ \langle\xi_t,w_t-w^{*}\rangle .
\tag{4}
\]

\noindent\textbf{[Step 3]}  Take expectation conditioned on $w_t$.  By Assumption~\ref{ass:noise}, $\E[\xi_t\,|\,w_t]=0$, thus
\[
\E\!\big[\langle\xi_t,w_t-w^{*}\rangle\,\big|\,w_t\big]=0.
\tag{5}
\]

\noindent\textbf{[Step 4]}  Apply convexity (Assumption~\ref{ass:convex}) to bound the deterministic part:
\[
f(w_t)-f(w^{*})\le \langle\nabla f(w_t), w_t-w^{*}\rangle .
\tag{6}
\]

\noindent\textbf{[Step 5]}  Combine (3)–(6) and take total expectation:
\[
\begin{aligned}
\E\!\big[\|w_{t+1}-w^{*}\|^{2}\big]
&\le \E\!\big[\|w_t-w^{*}\|^{2}\big]
-2\eta\,\E\!\big[f(w_t)-f(w^{*})\big]
+\eta^{2}\E\!\big[\|g_t\|^{2}\big].
\end{aligned}
\tag{7}
\]

\noindent\textbf{[Step 6]}  Expand $\|g_t\|^{2}$:
\[
\|g_t\|^{2}
= \|\nabla f(w_t)\|^{2}+2\langle\nabla f(w_t),\xi_t\rangle+\|\xi_t\|^{2}.
\tag{8}
\]

\noindent\textbf{[Step 7]}  Conditioned on $w_t$, the cross term vanishes:
\[
\E\!\big[\langle\nabla f(w_t),\xi_t\rangle\,\big|\,w_t\big]=0.
\tag{9}
\]

\noindent\textbf{[Step 8]}  Using the bounded‑gradient assumption (Assumption~\ref{ass:bound}) and the noise variance,
\[
\E\!\big[\|g_t\|^{2}\,\big|\,w_t\big]
\le G^{2}+\sigma^{2}d .
\tag{10}
\]

\noindent\textbf{[Step 9]}  Substitute (10) into (7) and remove the conditioning:
\[
\E\!\big[\|w_{t+1}-w^{*}\|^{2}\big]
\le \E\!\big[\|w_t-w^{*}\|^{2}\big]
-2\eta\,\E\!\big[f(w_t)-f(w^{*})\big]
+\eta^{2}\bigl(G^{2}+\sigma^{2}d\bigr).
\tag{11}
\]

\noindent\textbf{[Step 10]}  Rearrange (11) to isolate the expected suboptimality term:
\[
2\eta\,\E\!\big[f(w_t)-f(w^{*})\big]
\le \E\!\big[\|w_t-w^{*}\|^{2}\big]
-\E\!\big[\|w_{t+1}-w^{*}\|^{2}\big]
+\eta^{2}\bigl(G^{2}+\sigma^{2}d\bigr).
\tag{12}
\]

\noindent\textbf{[Step 11]}  Sum (12) over $t=0,\dots,T-1$; the telescoping sum yields
\[
2\eta\sum_{t=0}^{T-1}\E\!\big[f(w_t)-f(w^{*})\big]
\le \|w_0-w^{*}\|^{2}
+\eta^{2}T\bigl(G^{2}+\sigma^{2}d\bigr).
\tag{13}
\]

\noindent\textbf{[Step 12]}  Divide by $2\eta T$ and use convexity of $f$ to replace the average by the last iterate (or any convex combination).  For simplicity we keep the average:
\[
\frac{1}{T}\sum_{t=0}^{T-1}\E\!\big[f(w_t)-f(w^{*})\big]
\le
\frac{\|w_0-w^{*}\|^{2}}{2\eta T}
+\frac{\eta}{2}\bigl(G^{2}+\sigma^{2}d\bigr).
\tag{14}
\]

\noindent\textbf{[Step 13]}  Since $f$ is convex, the average objective is an upper bound on the expectation of the final iterate, i.e.
\[
\E\!\big[f(w_T)-f(w^{*})\big]\le
\frac{1}{T}\sum_{t=0}^{T-1}\E\!\big[f(w_t)-f(w^{*})\big].
\tag{15}
\]

\noindent\textbf{[Step 14]}  Combine (14) and (15) to obtain the claimed bound (2).  This completes the proof of Theorem~\ref{thm:conv}. \qed

%=====================================================================
\section*{Rate Derivation (With Constants)}
%=====================================================================
Choosing a constant stepsize
\[
\eta = \frac{\|w_0-w^{*}\|}{\sqrt{(G^{2}+\sigma^{2}d)T}}
\]
satisfies $\eta\le 1/L$ for sufficiently large $T$ (since $L$ is a fixed constant).  
Plugging this $\eta$ into (2) gives
\[
\E\!\big[f(w_T)-f(w^{*})\big]
\le
\frac{\|w_0-w^{*}\|}{\sqrt{T}}\sqrt{G^{2}+\sigma^{2}d},
\]
i.e. an $O(1/\sqrt{T})$ convergence rate with an explicit constant that grows with the noise scale $\sigma$ and the dimensionality $d$.

%=====================================================================
\section*{Special Cases}
%=====================================================================
\begin{itemize}
\item \textbf{No privacy noise ($\sigma=0$).}  The bound reduces to the classical SGD result
\[
\E[f(w_T)-f(w^{*})]\le \frac{\|w_0-w^{*}\|^{2}}{2\eta T}+\frac{\eta G^{2}}{2}.
\]

\item \textbf{Strongly convex $f$ (parameter $\mu>0$).}  Adding strong convexity yields a linear convergence term; the additional variance term $\eta\sigma^{2}d/2$ remains, leading to a steady‑state error proportional to $\sigma^{2}d/\mu$.

\item \textbf{High‑dimensional limit.}  When $d$ grows, the noise penalty scales linearly with $d$, emphasizing the need for dimension‑dependent clipping or advanced privacy mechanisms.
\end{itemize}

\end{document}